\begin{table}[htbp]
    \centering
    \caption[Descriptores de órbita y subespacios tangentes.]{Descriptores de las órbitas corregidas, muestreadas cada grado en los 25 parches fijos de validación. Cada cifra es una mediana entre parches. Las RMS de diferencias se promedian primero sobre los 360 ángulos; las razones respecto a RGB se calculan por parche antes de resumir. Las fracciones de energía armónica se refieren a la variación de la órbita latente completa.}
    \label{tab:descriptores-orbitas}
    \small
    \begin{tabular}{@{}lrr@{}}
        \toprule
        Descriptor & VAE base & VAE-\(\mathrm{SO}(2)\) \\
        \midrule
        RMS media de \(\Delta\mathbf z\) & 0,406 & 0,386 \\
        RMS media de \(\Delta^2\mathbf z\) & 0,384 & 0,339 \\
        Longitud RMS acumulada & 146,223 & 138,989 \\
        Longitud relativa a la órbita RGB & 3,558 & 3,405 \\
        Linealidad local relativa a RGB & 0,649 & 0,601 \\
        Curvatura mediana relativa a RGB & 0,106 & 0,098 \\
        Fracción de energía en armónicos 1--6 & 0,107 & 0,106 \\
        Número efectivo de frecuencias & 4,135 & 3,982 \\
        Fracción de variación capturada por PC1--PC6 & 0,111 & 0,109 \\
        Direcciones tangentes para el 90\% & 120 & 114 \\
        Direcciones tangentes para el 95\% & 142 & 133 \\
        Direcciones tangentes para el 99\% & 211 & 176 \\
        Fracción de energía tangente local en una dirección & 0,458 & 0,467 \\
        \bottomrule
    \end{tabular}
\end{table}
